\chapter{Discussion}
\label{ch:discussion}

\section{What the equivalence invariant bought}

The decision that shaped everything else was to fix the reduction chunk grid by
problem size alone. It looked at first like a small implementation detail and it
turned out to be the foundation of the whole comparison.

Because every shared memory backend produces bit identical iterates, the timing
differences between them cannot be attributed to any numerical difference: there
is none. That removes an entire class of objection from the results chapter. It
also turned the test suite into something sharper than a regression check.
Asserting exact equality means a test fails on the first bit that differs, which
is a much more sensitive detector than a tolerance that would have absorbed real
faults silently. The fused multiply add discrepancy in Section
\ref{sec:cuda-design} was one bit in the last place, and a tolerance based test
would never have reported it.

The cost is real and is measured in Section \ref{sec:reduction-cost} rather than
excused. It is worth paying here because reproducibility is a deliverable of
this project; a production solver with different priorities might reasonably
choose the native reduction and say so.

\section{The methods that do not parallelise, and why they are still here}

Natural ordering Gauss Seidel, SOR and block Gauss Seidel are sequentially
dependent. There was a real temptation to substitute a red black reordering
whenever more than one worker was present, which would have produced attractive
scaling curves for all of them.

That would have been a lie of a particular kind: not a wrong number, but a right
number for a different algorithm. The library instead preserves exact sequential
semantics on every backend, including across MPI ranks through a pipelined token
chain. The consequence is that these methods appear in the scaling figures as
essentially flat lines, and that is the correct result. The reader learns
something true, which is that a method with a sequential recurrence has nothing
to offer a parallel machine, and that red black ordering is the price one pays
to get parallelism back.

Charging that price explicitly also improved the device comparison. Because the
red black penalty is measured (\pnlRedBlackIterations{} against
\pnlNaturalIterations{} iterations on a \pnlRedBlackGrid{} grid,
\pnlRedBlackPenalty{} percent), it can be applied to both the CPU and the GPU
side, and the device figure is not silently inflated by a change of algorithm.

\section{Where the environment set the limits}

Three findings came from the environment rather than from the mathematics, and
all three changed what could be claimed.

\paragraph{The hybrid core topology is invisible from inside WSL2.} The guest
synthesises a homogeneous machine, and thread affinity binds to a virtual
processor the hypervisor may place anywhere. The response was to measure rather
than assume, and to take the knee from the aggregate scaling curve, which does
not depend on identifying individual cores. That is a weaker claim than the
objective originally envisaged and it is the strongest one the environment
supports.

\paragraph{The CUDA toolchain and the host toolchain cannot be the same.}
CUDA 13.3 rejects any host compiler newer than GCC 15 and cannot parse GCC 15's
own headers either, so no released GCC serves as both. Resolving this through a
C ABI boundary was not a workaround so much as a clarification: the boundary was
already required by the rule that backends do not leak their model's types, and
enforcing it at the ABI level made that rule impossible to violate by accident.

\paragraph{MPI version macros describe conformance, not capability.} OpenMPI
5.0.10 advertises MPI 3.1 while implementing many MPI 4.0 entry points. Guarding
on the version macro would have disabled working functionality; assuming the
functions were present would have broken elsewhere. Probing for symbols is the
only approach that is both correct and portable.

\section{Things that were wrong before they were right}

Four faults are worth naming because each was the kind that produces a plausible
number rather than an obvious failure.

\paragraph{The manufactured solution was an eigenvector.} With
$u = \sin(\pi x)\sin(\pi y)$ the Krylov subspace is one dimensional and
conjugate gradient converges in one iteration at every grid size. That looks
like a spectacular result and is a degenerate problem. It was caught because a
test asserted the iteration count should \emph{grow} with the grid, not merely
that the method converged. A test that had only checked convergence would have
passed forever.

\paragraph{A default that was right for one solver and wrong for another.} The
relaxation factor defaulted to one, which is harmless for SOR, where it means
Gauss Seidel, and catastrophic for Richardson, where on the Poisson operator it
gives an iteration matrix with spectral radius seven. The lesson generalises: a
shared options struct wants defaults that mean "ask the component", not defaults
that happen to suit whichever component was written first.

\paragraph{The resumable sweep was not resumable.} The skip test sat after the
work rather than before it, so a restart re-ran every configuration and skipped
only the write. The counters reported plausible numbers throughout, which is why
it survived until a restart was timed rather than merely observed.

\paragraph{A tolerance below the arithmetic.} Two tests asked for a relative
residual that stationary iterations on this operator cannot reach, because over
relaxation raises their rounding limited floor above it. The failure appeared
only when the solver those tests used changed, which made it look like a
regression in the solver rather than an impossible request. Measuring the floor
turned a confusing symptom into a table, and the tests now ask for a tolerance
above it.

\section{What the traffic model changes, and what it does not}
\label{sec:traffic-correction}

The open question of Section \ref{sec:preregistered-traffic} is whether a Jacobi
pass over the five point stencil moves 24 bytes per unknown or 32, the extra
eight being the output line an ordinary store has to fetch before it can
overwrite it. It is worth being exact about what turns on the answer, because
less turns on it than an earlier draft of this section claimed.

An efficiency is one achieved bandwidth over another, and it says something only
when both sides are counted the same way. The triad this report divides by is a
plain C++ loop, so its own store pays exactly the charge the counted model
charges the kernel, and under the counted model the denominator is recounted
too. The charge then cancels, and a counted efficiency is the declared
efficiency times the row's pass count over its sweep count. The byte model moves
the achieved bandwidth figure in GiB/s and leaves the efficiency where it was,
on the host and on the device alike. Host Jacobi stands at
\pnlJacobiHostEfficiency{} percent of the plain triad and the device at
\pnlJacobiDeviceEfficiency{} percent of its own, under either model, and the two
percentage columns of Table \ref{tab:device-comparison} agree on those rows to
the digit they print.

The Jacobi row is therefore not an exception waiting for a byte model to remove
it, and it is not an exception at all. Its device over host efficiency ratio is
\pnlJacobiEfficiencyRatio{} in this generation, where release 1.0.0's published
generation put it far higher; the before and after per phase is in
\texttt{PROGRESS.md} rather than typed here. What moved it is not the traffic
model. It is that the per iteration full state copy is gone from the Jacobi path
and that the timed region no longer holds the allocation, so the host figure is
a measurement of the sweep rather than of the sweep plus a serial copy plus a
workspace. That is the correction that mattered, and it is a correction to the
measurement and not to the model that reads it.

The other three kernels gave that conclusion independently and Jacobi now joins
them. At \pnlDeviceSize{} unknowns the device over host efficiency ratio is
\pnlRedBlackEfficiencyRatio{} for red black Gauss Seidel and
\pnlCgEfficiencyRatio{} for conjugate gradient, against Jacobi's
\pnlJacobiEfficiencyRatio{}, so on these kernels the great majority of the
observed speedup is the memory system and little of it is attributable to the
port. That holds under the conservative count, without the correction the
traffic model question was about, which is a better outcome than needing it.

\paragraph{The sentence this replaces.} An earlier draft of this section said
that under read for ownership the host achieved bandwidth and the host
efficiency both rise by a third while the device figures do not move, and that
host and device efficiency converge on the Jacobi row as a result. The first
half is right about GiB/s and wrong about the efficiency, and the rest was the
same mismatch: it divided a numerator counted one way by a denominator counted
the other. The pre registered sentence quoted in Section
\ref{sec:preregistered-traffic} carries the same reasoning, and it is quoted
there exactly as it was written, because a pre registration edited after the
measurement is not one. This paragraph is the correction to it, and the outcome
that sentence was written for is not the one the measurement selected.

The report does not choose between the two models on its own authority. The rule
was fixed before the measurement, the measurement is the publication session's,
and the generator applies the rule and writes the sentence, which Section
\ref{sec:device-results} carries. That sentence records the outcome as
\pnlTrafficModel{}, which is this instrument saying it cannot separate the two
counts on this machine rather than choosing between them.

Both counts are therefore carried in every table and every figure, every claim
in this chapter that would change under the other model is stated conditionally,
and no conclusion here rests on one of them in a way that would fail under the
other. For a method whose pass count exceeds its sweep count the two columns
bracket the traffic rather than compete to describe it, which Section
\ref{sec:device-results} sets out; for Jacobi they are the same number.

\section{Related work}
\label{sec:related-work}

Deterministic reduction is established technique and nothing here invented it.
Saying so is not modesty: the contribution of this project is a narrower thing
than the reduction, and not naming the prior art both reads as unawareness of
the field and undersells the work.

\paragraph{Reproducible arithmetic.} Intel's Math Kernel Library offers
Conditional Numerical Reproducibility \cite{mklcnr}, which gives bitwise
identical results across runs and across a range of processors subject to
conditions the vendor documents: a fixed thread count, aligned arrays, and a
selected code path. It states the shape of the problem clearly, which is that
reproducibility is bought by fixing what a performance library would otherwise
be free to vary. Demmel and Nguyen give summation algorithms whose result does
not depend on the order or the number of summands, at bounded cost and with an
error bound \cite{demmel2013,demmel2015}, and Ahrens, Demmel and Nguyen develop
the binning algorithm behind ReproBLAS \cite{ahrens2020}. That is a stronger
guarantee than the one made here: ReproBLAS is reproducible across different
decompositions, whereas the fixed chunk grid of this library is reproducible
because the decomposition is held fixed.

\paragraph{Performance portability.} Kokkos
\cite{edwards2014,trott2022} and RAJA \cite{beckingsale2019} are the established
C++ layers for writing a kernel once and running it over several execution
models including GPUs. Both provide reductions and both document what is and is
not deterministic about them. They solve the portability problem this library
also solves, at a scale and maturity this library does not approach.

\paragraph{Solver libraries.} Ginkgo \cite{anzt2022} is a modern C++ sparse
linear algebra framework built around linear operators, with GPU backends and a
strong benchmarking culture. PETSc \cite{balay1997,petsc2026} is the reference
distributed solver toolkit, whose Krylov methods and preconditioners this
repository implements a small subset of. hypre \cite{falgout2002} is the
reference library for scalable multigrid preconditioners. Each is a real answer
to the question of what to use to solve a system, and this library is not.

\paragraph{What this project adds.} Not the deterministic reduction. Three
things, and they are narrower.

\begin{enumerate}
  \item \textbf{One numerical invariant held across seven execution models at
        once}, a GPU and a distributed backend among them. The comparison
        between programming models in Chapter \ref{ch:results} is only
        meaningful because the thing being executed is provably identical, and
        holding one contract across that particular set of models is the part
        that is unusual.
  \item \textbf{Asserted as exact equality in a test suite}, rather than
        described in documentation. The equivalence suite compares with
        \texttt{==} across twelve solvers, two problem families, four backends
        and seven worker counts, and fails on the first bit that differs. The
        one bit contraction discrepancy of Section \ref{sec:cuda-design} was
        found that way and would have vanished into any tolerance.
  \item \textbf{Priced.} What the invariant costs is measured against each
        model's own native reduction and reported with its run to run spread,
        including where that spread is wide enough that the honest answer is
        that the cost cannot be separated from the noise.
\end{enumerate}

\paragraph{What is missing, and when it arrives.} There is no baseline
comparison in this report. Nothing here says what a competent existing
implementation achieves on this machine on this problem, so a statement of the
form "this solver reaches some percentage of host STREAM" cannot be read as good
or bad. That is a real gap and it is stated as one. A single baseline against
PETSc is planned for release 1.2.0: the same stencil, \texttt{KSPCG} with
\texttt{PCJACOBI} and with \texttt{PCSOR}, the same tolerance, the same machine
and the same reported metrics, emitting rows in the schema the sweep already
uses. One baseline is enough, and the result is worth publishing whichever way
it comes out. \texttt{docs/RELATED\_WORK.md} carries this section in longer form.

\section{What a reader should take away}

For someone choosing a programming model, the results chapter says what each
model costs on identical work, with the numerical variable removed by
construction. For someone choosing a method, the convergence tables say what each
one costs in iterations, independent of any machine. For someone weighing a GPU,
Section \ref{sec:device-results} separates the part of the speedup that is the
memory system from the part that is the implementation, which is the separation
that determines whether the speedup would transfer to their problem.

The most transferable conclusion is methodological rather than numerical. A
comparison between execution models is only meaningful if the thing being
executed is identical, and the only way to know that is to test it at the bit
level. Everything else in this project follows from taking that seriously.
